\section{Особенности реализации}
Программа написана в виде единого скрипта на языке C++. Реализовано текстовое меню,
с помощью которого пользователь может выбрать направление перевода:
из десятичной системы счисления в двоичный код в формате bf16 или из
двоичного кода bf16 --- в десятичную систему счисления,
также можно выйти из программы.
Реализована обработка некорректного
пользовательского ввода: если вместо номера пункта меню введено число
вне диапазона от 0 до 3 или введены не цифры, программа выводит
сообщение об ошибке и просит ввести значение заново; также проверяется
корректность вводимого числа для перевода (десятичное число) и
вводимого кода (строка ровно из 16 символов, состоящая только из 0 и 1).

\subsection{Основные функции}

Функция \textbf{binaryToIntPart} предназначена для перевода целого числа
из двоичной системы счисления в десятичную.

\begin{verbatim}
Алгоритм binaryToIntPart(binaryNumber).
Вход: binaryNumber --- строка, двоичное число.
Выход: строка --- десятичное представление числа.

начало
    result = 0

    если binaryNumber равно "0", то
        вернуть "0"
    конец если

    n = длина(binaryNumber)

    для i от n - 1 до 0 (шаг -1) выполнять
        если binaryNumber[i] равно '1', то
            result = result + 2 в степени (n - 1 - i)
        конец если
    конец для

    вернуть result, переведённое в строку
конец
\end{verbatim}

Функция \textbf{intPartToBinary} предназначена для
перевода целого числа из десятичной системы счисления
в двоичную.

\begin{verbatim}
Алгоритм intPartToBinary(decimalNumber).
Вход: decimalNumber --- целое число.
Выход: строка --- двоичное представление числа.

начало
    если decimalNumber равно 0, то
        вернуть "0"
    конец если

    result = пустая строка

    пока decimalNumber больше или равно 1 выполнять
        если decimalNumber нечётное, то
            добавить '1' в конец result
        иначе
            добавить '0' в конец result
        конец если
        decimalNumber = decimalNumber / 2 (целочисленно)
    конец пока

    развернуть result
    вернуть result
конец
\end{verbatim}

Функция \textbf{fractionToBinary} предназначена для перевода
дробной части числа из десятичной системы
счисления в двоичную с фиксированной точностью в 10 знаков.

\begin{verbatim}
Алгоритм fractionToBinary(fraction).
Вход: fraction --- дробная часть числа (0 <= fraction < 1).
Выход: строка из 10 символов --- двоичное представление 
дробной части.

начало
    result = пустая строка

    повторить 10 раз
        fraction = fraction * 2
        если fraction больше или равно 1, то
            добавить '1' в конец result
            fraction = fraction - 1
        иначе
            добавить '0' в конец result
        конец если
    конец повторения

    вернуть result
конец
\end{verbatim}

Функция \textbf{decimalTo\_bf16} предназначена для перевода вещественного
числа из десятичной системы счисления в двоичный код в формате bf16.

\begin{verbatim}
Алгоритм decimalTo_bf16(decimalNumber).
Вход: decimalNumber --- вещественное число.
Выход: строка из 16 символов --- код числа в формате bf16.

result = пустая строка

начало
    // знаковый бит
    если decimalNumber меньше 0, то
        result = result + '1'
    иначе
        result = result + '0'
    конец если

    absVal = модуль(decimalNumber)

    // случай нулевого числа
    если absVal равно 0, то
        вернуть result + "00000000" + "0000000"
    конец если

    // поиск несмещённой экспоненты
    если absVal больше или равно 1, то
        wholePart = целая часть(absVal)
        exponent = длина(двоичная_запись(wholePart)) - 1
    иначе
        exponent = 0
        copy = absVal
        пока copy меньше 1 выполнять
            copy = copy * 2
            exponent = exponent - 1
        конец пока
    конец если

    // случай переполнения
    если exponent + 127 больше или равно 255, то
        вернуть result + "11111111" + "0000000"
    конец если

    // денормализованное число
    если exponent меньше -126, то
        scaled = absVal / (2 в степени -126)
        mantissaBits = первые 7 бит от fractionToBinary(scaled)
        вернуть result + "00000000" + mantissaBits
    конец если

    // нормализованное число
    если absVal больше или равно 1, то
        wholePart = целая часть(absVal)
        fractionPart = absVal - wholePart
        mantissaBits = первые 7 бит от
            (двоичная_запись(wholePart) без первого символа + 
            fractionToBinary(fractionPart))
    иначе
        normalizedFraction = absVal / (2 в степени exponent) - 1
        mantissaBits = первые 7 бит от 
        fractionToBinary(normalizedFraction)
    конец если

    exponentBinary = двоичная_запись(exponent + 127),
    дополненная нулями слева до 8 символов

    result = result + exponentBinary + mantissaBits
    вернуть result
конец
\end{verbatim}

Функция \textbf{binaryBF16ToDecimal} предназначена для перевода числа,
представленного в формате bf16, из двоичной системы счисления в десятичную.

\begin{verbatim}
Алгоритм binaryBF16ToDecimal(bits).
Вход: bits --- строка из 16 символов, код числа в формате bf16.
Выход: вещественное число --- десятичное представление.

начало
    signBit = первый символ bits
    exponentBits = символы с 2 по 9 (8 символов)
    mantissaBits = символы с 10 по 16 (7 символов)

    E = десятичное представление(exponentBits)
    M = десятичное представление(mantissaBits)

    если E равно 0 и M равно 0, то
        value = 0
    иначе если E равно 0 и M не равно 0, то
        value = (M / 128) * (2 в степени -126)
    иначе если E равно 255, то
        если M равно 0, то
            value = бесконечность 
        иначе
            value = неопределённость 
    иначе
        unbiasedExponent = E - 127
        value = (M / 128 + 1) * (2 в степени unbiasedExponent)
    конец если

    если signBit равен '1', то
        value = -value
    конец если

    вернуть value
конец
\end{verbatim}